\section{结语：帝国形成取决于接收的尺度}

蒙古的胜利把原先相互竞争的辽、宋、西夏、金和南宋空间置于新的权力关系中，却没有使区域差异消失。站赤、军政组织、行省和税源调度依赖既有道路、城市、人员与地方知识；同一制度在河西、华北、江南和海港的运作也不可能相同。征服扩大了流动和交换的范围，同时把转运、劳役与暴力的成本带入更广泛的社会。

因此，元朝不是征服叙事的自然终点，而是接收问题的下一阶段。大都与江南、草原与海域之间的联系，需要以财政、文书、宗教机构和多语中介反复维持。下一章将考察这些联系怎样形成一个不均衡的行省帝国，并检验中枢制度在不同地方实际能够做到什么。
